\section{Turbulence Onset as a Regime I \texorpdfstring{$\to$}{to} II Transition}
\label{sec:turbulence_onset}

\begin{objectivebox}
\begin{itemize}
    \item Model turbulence onset as a Regime~I $\to$ II transition where the critical Reynolds number maps to the saturation threshold.
\end{itemize}
\end{objectivebox}

This section treats the Regime~I $\to$ II onset transition. The Regime~III boundary and the resolution of the Kolmogorov cascade singularity are covered in Volume~3, Chapter~16 (The Kolmogorov Spectral Cutoff).

\subsection{Hypothesis}

The laminar-to-turbulent transition in fluid flow occurs at a
critical Reynolds number $\mathrm{Re}_c$ that depends on geometry.
The AVE framework identifies this transition as the Regime~I $\to$
Regime~II boundary crossing, where the driving parameter (flow
velocity, or equivalently the viscous stress) exceeds the first
universal regime boundary $r_1 = \sqrt{2\alpha}$ of the canonical
Four-Regime Map (\kbleaf{vol1/operators-and-regimes/ch7-regime-map/four-regimes.md}).

The Reynolds number maps to a normalised drive parameter:
\begin{equation}
  r = \frac{\mathrm{Re}}{\mathrm{Re}_{\max}},
\end{equation}
where $\mathrm{Re}_{\max}$ is the geometry-dependent maximum
stable Reynolds number.  The universal prediction is that the
onset of turbulence (first appearance of self-sustaining vortical
structures) occurs at $r = r_1 = \sqrt{2\alpha} \approx 0.1208$.

\subsection{Mapping to Fluid Variables}

The regime boundary $r_1$ sets the ratio of inertial driving to
viscous dissipation at which the laminar standing wave (Poiseuille
flow) first departs the linear regime and begins to shed energy
into higher modes (turbulent eddies).  The saturation operator
governs the transition:
\begin{equation}
  S(r) = \sqrt{1 - r^2}, \qquad r = \frac{\mathrm{Re}}{\mathrm{Re}_{\max}}.
\end{equation}
This is the bare quarter-arc kernel ($f_{\text{yield}} = 1$) of the
canonical Four-Regime Map; the structural-fluid yield prefactor
$f_{\text{yield}}$ applies to the macroscopic rupture avalanche, not
to the onset mapping.
For $r < r_1 = \sqrt{2\alpha}$, the flow sits in Regime~I with
$S > 0.993$: the saturation correction is sub-$\alpha$ and the flow
remains laminar.  At the critical Reynolds number
$\mathrm{Re}_c = r_1 \, \mathrm{Re}_{\max}$, the drive crosses the
Regime~I $\to$ II boundary at $S(r_1) = \sqrt{1 - 2\alpha} \approx 0.993$,
where the laminar mode first leaves the linear regime and turbulent
bursts appear.  Full saturation ($S = 0$) is the Regime~IV rupture
endpoint at $r = 1$, not the onset.

\subsection{Testable Prediction}

The intermittency fraction (fraction of time the flow exhibits
turbulent bursts) near the critical Reynolds number should follow
$1 - S(r)$ with $r = \mathrm{Re}/\mathrm{Re}_{\max}$, measurable via
hot-wire anemometry in pipe flow at controlled Reynolds numbers.

\begin{summarybox}
\begin{itemize}
    \item Turbulence onset is modelled as a Regime~I $\to$ II transition: the critical Reynolds number $\mathrm{Re}_c = \sqrt{2\alpha}\,\mathrm{Re}_{\max}$ maps to the saturation threshold $S(r)$ with $r = \mathrm{Re}/\mathrm{Re}_{\max}$, predicting intermittent turbulent bursts near the transition boundary.
\end{itemize}
\end{summarybox}

\begin{exercisebox}
\begin{enumerate}
    \item Design a hot-wire anemometry experiment to verify that turbulent burst statistics near $\mathrm{Re}_c$ follow the saturation function $1 - S(r)$ with $r = \mathrm{Re}/\mathrm{Re}_{\max}$.
\end{enumerate}
\end{exercisebox}

